\chapter{Comunicação Serial e Depuração}
\label{cap:serial}

Sem teclado, tela ou depurador visual conectado por padrão, a \textbf{comunicação serial} é a principal janela do desenvolvedor para dentro de um firmware em execução — é como saber o que o programa está ``pensando'' a cada momento, essencial tanto para aprender quanto para depurar problemas.

\section{UART e o monitor serial}

O \ESP{} se comunica com o computador através de uma interface \textbf{UART} (\emph{Universal Asynchronous Receiver/Transmitter}), tipicamente ligada a um chip conversor USB-serial na própria placa de desenvolvimento. No framework Arduino, essa comunicação é controlada pelo objeto \code{Serial}, que precisa ser iniciado uma única vez, em \code{setup()}, com a velocidade de comunicação desejada.

\begin{codigoc}{src/main.cpp — contador via Serial}
void setup() {
    Serial.begin(115200);
}

void loop() {
    static int contador = 0;

    Serial.print("Contador: ");
    Serial.println(contador);
    contador++;

    delay(1000);
}
\end{codigoc}

Para visualizar essa saída, use o \textbf{Serial Monitor} do PlatformIO (\cref{cap:platformio}) — ele abre uma conexão serial com a placa na velocidade definida por \code{monitor\_speed} no \code{platformio.ini} (que deve coincidir com o valor passado a \code{Serial.begin}) e exibe, em tempo real, tudo o que o firmware escreve.

\subsection{As variantes de \code{Serial.print}}

\begin{table}[h]
\centering
\begin{tabular}{@{}ll@{}}
\toprule
\textbf{Função} & \textbf{Comportamento} \\
\midrule
\code{Serial.print(x)}     & imprime \code{x}, sem quebra de linha ao final \\
\code{Serial.println(x)}   & imprime \code{x}, seguido de uma quebra de linha \\
\code{Serial.printf(fmt, \ldots)} & imprime com especificadores de formato, como \code{printf} \\
\bottomrule
\end{tabular}
\caption{Formas de escrever na porta serial no framework Arduino.}
\end{table}

\begin{dica}
\code{Serial.printf} aceita exatamente os mesmos especificadores de formato apresentados no \cref{cap:entrada-saida} (\code{\%d}, \code{\%f}, \code{\%s}, \code{\%.2f}, e assim por diante) — todo aquele conhecimento é diretamente reaproveitável aqui, bastando trocar \code{printf(...)} por \code{Serial.printf(...)}.
\end{dica}

\section{Saída como ferramenta de depuração, agora no firmware}

A técnica de \emph{print debugging} apresentada no \cref{sec:printf-debug} — imprimir o valor de uma variável em pontos estratégicos do código para entender o comportamento de um programa — se transporta quase sem alterações para o firmware, usando \code{Serial} no lugar de \code{printf}:

\begin{codigoc}{Depurando uma leitura de sensor}
int leitura_bruta = analogRead(PINO_SENSOR);
Serial.print("[DEBUG] leitura_bruta = ");
Serial.println(leitura_bruta);

float tensao = (leitura_bruta / 4095.0f) * 3.3f;
Serial.print("[DEBUG] tensao = ");
Serial.println(tensao);
\end{codigoc}

\section{Log por nível de severidade}

O núcleo Arduino do \ESP{} também oferece um sistema de \textbf{log por nível de severidade} — uma alternativa mais expressiva que \code{Serial.print} puro, com macros que já indicam a gravidade da mensagem:

\begin{codigoc}{Usando as macros de log}
void setup() {
    Serial.begin(115200);
}

void loop() {
    int leitura = analogRead(34);

    log_i("Sistema em operacao normal.");
    if (leitura < 100) {
        log_w("Leitura muito baixa: %d", leitura);
    }
    if (leitura == 0) {
        log_e("Sensor parece desconectado!");
    }

    delay(1000);
}
\end{codigoc}

\begin{table}[h]
\centering
\begin{tabular}{@{}lll@{}}
\toprule
\textbf{Macro} & \textbf{Nível} & \textbf{Uso típico} \\
\midrule
\code{log\_e} & Error   & falhas que impedem o funcionamento normal \\
\code{log\_w} & Warning & situações anormais, mas não fatais \\
\code{log\_i} & Info    & mensagens informativas de operação normal \\
\code{log\_d} & Debug   & detalhes úteis apenas durante desenvolvimento \\
\code{log\_v} & Verbose & o nível mais detalhado, quase sempre desativado \\
\bottomrule
\end{tabular}
\caption{Níveis de log disponíveis no núcleo Arduino do ESP32.}
\end{table}

O nível mínimo de mensagens exibidas é controlado por uma opção de compilação, ajustável diretamente no \code{platformio.ini}:

\begin{codigoini}{platformio.ini com nível de log elevado}
[env:esp32dev]
platform = espressif32
board = esp32dev
framework = arduino
monitor_speed = 115200
build_flags = -DCORE_DEBUG_LEVEL=4   ; 0=Nenhum ... 5=Verbose
\end{codigoini}

\begin{dica}
A grande vantagem de usar \code{log\_x} em vez de \code{Serial.print} puro é justamente essa: o nível mínimo de log pode ser reduzido em uma versão ``de produção'' do firmware (silenciando mensagens de depuração detalhadas) sem precisar remover ou comentar cada chamada manualmente — basta alterar um único número em \code{platformio.ini}.
\end{dica}

\begin{esp32box}
Uma armadilha clássica ao depurar por log em sistemas embarcados: mensagens de log \emph{também} consomem tempo de processamento e, no caso da UART, banda de comunicação limitada. Excesso de logs dentro de um laço executado milhares de vezes por segundo pode, ele próprio, atrasar o firmware o suficiente para mascarar ou até causar o problema que se está tentando diagnosticar — um efeito colateral que vale ter em mente ao decidir onde e com que frequência logar.
\end{esp32box}
